% ---- pg-293 （书页 281） ----

也是随机变量), 其行为由一组关于分布的假定来表征.

具体地说, 单个响应的线性回归模型取如下形式:
\[
Y = \beta_0+\beta_1 z_1+\cdots+\beta_r z_r+\epsilon
\]
\[
(\text{响应})=[\text{均值}(\text{依赖于 } z_1,\cdots,z_r)]+(\text{误差})
\]
线性一词指均值为未知参数 $\beta_0,\beta_1,\cdots,\beta_r$ 的线性函数, 预测变量在模型中不一定是一阶项.

如将模型用 $Y$ 的 $n$ 个观测值和与之相联系的诸 $z_i$ 值来表示时, 则整个模型变为
\phantomsection\label{eq:7-1}
\begin{equation}
\begin{aligned}
Y_1 &= \beta_0+\beta_1 z_{11}+\beta_2 z_{12}+\cdots+\beta_r z_{1r}+\epsilon_1\\
Y_2 &= \beta_0+\beta_1 z_{21}+\beta_2 z_{22}+\cdots+\beta_r z_{2r}+\epsilon_2\\
&\quad\vdots\\
Y_n &= \beta_0+\beta_1 z_{n1}+\beta_2 z_{n2}+\cdots+\beta_r z_{nr}+\epsilon_n
\end{aligned}
\tag{7-1}
\end{equation}
其中误差项假定具有以下性质:
\phantomsection\label{eq:7-2}
\begin{equation}
\begin{array}{ll}
1. & E(\epsilon_j)=0;\\
2. & \text{Var}(\epsilon_j)=\sigma^2\ (\text{常数});\\
3. & \text{Cov}(\epsilon_j,\epsilon_k)=0,\ j\neq k.
\end{array}
\tag{7-2}
\end{equation}
用矩阵形式表示, \hyperref[eq:7-1]{式(7-1)} 变成
\[
\begin{bmatrix}Y_1\\Y_2\\\vdots\\Y_n\end{bmatrix}
=
\begin{bmatrix}
1&z_{11}&z_{12}&\cdots&z_{1r}\\
1&z_{21}&z_{22}&\cdots&z_{2r}\\
\vdots&\vdots&\vdots&&\vdots\\
1&z_{n1}&z_{n2}&\cdots&z_{nr}
\end{bmatrix}
\begin{bmatrix}\beta_0\\\beta_1\\\vdots\\\beta_r\end{bmatrix}
+
\begin{bmatrix}\epsilon_1\\\epsilon_2\\\vdots\\\epsilon_n\end{bmatrix}
\]
或
\[
\underset{(n\times 1)}{\mathbf{Y}}
=
\underset{(n\times(r+1))}{\mathbf{Z}}\;
\underset{((r+1)\times 1)}{\boldsymbol{\beta}}
+
\underset{(n\times 1)}{\boldsymbol{\epsilon}}
\]
而\hyperref[eq:7-2]{式(7-2)} 中的假定可表示为:
\[
\begin{array}{ll}
1. & E(\boldsymbol{\epsilon})=\mathbf{0};\\
2. & \text{Cov}(\boldsymbol{\epsilon})=E(\boldsymbol{\epsilon}\boldsymbol{\epsilon}')=\sigma^2\mathbf{I}.
\end{array}
\]
注意设计矩阵 $\mathbf{Z}$ 的第一列为常数项 $\beta_0$ 的系数. 习惯上引入一个人工变量 $z_{j0}=1$, 使得
\[
\beta_0+\beta_1 z_{j1}+\cdots+\beta_r z_{jr}=\beta_0 z_{j0}+\beta_1 z_{j1}+\cdots+\beta_r z_{jr}
\]
$\mathbf{Z}$ 的每一列由相应预测变量的 $n$ 个值组成, 而 $\mathbf{Z}$ 的第 $j$ 行则包含第 $j$ 次试验中所有预测变量的值.

\begin{mvframe}[经典线性回归模型]
\[
\underset{(n\times 1)}{\mathbf{Y}}
=
\underset{(n\times(r+1))}{\mathbf{Z}}\;
\underset{((r+1)\times 1)}{\boldsymbol{\beta}}
+
\underset{(n\times 1)}{\boldsymbol{\epsilon}},
\]
\phantomsection\label{eq:7-3}
\begin{equation}
E(\boldsymbol{\epsilon})=\underset{(n\times 1)}{\mathbf{0}},\qquad
\text{Cov}(\boldsymbol{\epsilon})=\underset{(n\times n)}{\sigma^2\mathbf{I}},
\tag{7-3}
\end{equation}
\end{mvframe}

其中 $\boldsymbol{\beta}$ 和 $\sigma^2$ 为未知参数, 设计矩阵 $\mathbf{Z}$ 的第 $j$ 行为 $[z_{j0},z_{j1},\cdots,z_{jr}]$.

虽然\hyperref[eq:7-2]{式(7-2)} 中关于误差项的假定很宽松, 但稍后我们将增添一个正态性假定, 以便能构造置信区间和进行假设检验.

下面提供一些有关线性回归模型的例子.

\medskip
\noindent\textbf{例 7.1}\quad (拟合直线回归模型)

用数据

\begin{mvtable}
\begin{tabular}{c|ccccc}
\hline
$z_1$&0&1&2&3&4\\
\hline
$y$&1&4&3&8&9\\
\hline
\end{tabular}
\end{mvtable}
